% --- Archon physics formalization source begin ---
% archon:physics
% archon:covers PhyXMiniProblems/problem_phyx_mini_0818.lean
% archon:source-report reports/phyx_mini/problem_phyx_mini_0818.source.json

\chapter{Physics problem phyx\_mini\_0818}
\label{ch:PhyXMiniProblems_problem_phyx_mini_0818}

\paragraph{Problem source.}
\#\# Physical scenario

You are designing an airplane propeller that is to turn at 2400 rpm as shown in figure. The forward airspeed of the plane is to be \textbackslash{}( 75.0 \textbackslash{}text\{ m/s\} \textbackslash{}), and the speed of the propeller tips through the air must not exceed \textbackslash{}( 270 \textbackslash{}text\{ m/s\} \textbackslash{}). (This is about \textbackslash{}( 80\textbackslash{}\% \textbackslash{}) of the speed of sound in air. If the propeller tips moved faster, they would produce a lot of noise.)

\#\# Question

What is the maximum possible propeller radius?

\#\# Answer choices

- A: A: 1.87m

- B: B: 1.03m

- C: C: 2.12m

- D: D: 2.08m

\#\# Figure caption (auxiliary; use the image as primary evidence)

The image is an illustration of an airplane with a propeller. The airplane is moving to the right, and the propeller is positioned at the front. Several annotations and arrows describe the motion and speed of the airplane and propeller:

1. **Airplane and Propeller:**
   - The airplane is depicted with a cockpit and wings.
   - The propeller, shown at the nose of the airplane, is colored orange.

2. **Arrows and Annotations:**
   - A green arrow labeled \textbackslash{}( v\_\{\textbackslash{}text\{plane\}\} = 75.0 \textbackslash{}, \textbackslash{}text\{m/s\} \textbackslash{}) points to the right, indicating the velocity of the airplane.
   - A black arrow shows a rotational direction around the propeller, marked with \textbackslash{}( 2400 \textbackslash{}, \textbackslash{}text\{rev/min\} \textbackslash{}), suggesting the rotation speed of the propeller.
   - A green downward arrow labeled \textbackslash{}( v\_\{\textbackslash{}text\{tan\}\} = r\textbackslash{}omega \textbackslash{}) indicates the tangential velocity of the propeller blades.
   - A horizontal black arrow from the center of the propeller, labeled \textbackslash{}( r \textbackslash{}), shows the radius from the center to the tip of the propeller blade.

This diagram seems to illustrate concepts of rotational motion and tangential velocity in relation to the functionality of a propeller on an airplane.

\#\# Dataset metadata

Rotational Motion; Physical Model Grounding Reasoning, Multi-Formula Reasoning

\paragraph{Recorded answer/context.}
B: B: 1.03m

\paragraph{Figure/image path.}
/root/proposal\_for\_physic/science-mango/phyx\_mini\_run/phyx\_data/test\_image/818.png

\paragraph{Formalization target.}
create a compiling Lean file with sorry bodies at `PhyXMiniProblems/problem\_phyx\_mini\_0818.lean`. The Lean declarations must preserve the physical quantities, dimensions or dimensional roles, figure labels, governing-law hypotheses, and final relation expressed by this problem.
Use Mathlib/Physlib names found through LeanExplore where available. If a physics API is missing, introduce faithful local abstractions rather than scalar placeholder aliases.

\begin{theorem}[Physics formalization target]
\label{thm:physics:phyx_mini_0818:target}
\lean{PhyXMiniProblems.ProblemPhyXMini0818.maximumPossiblePropellerRadius_is_one_point_zero_three_meters}
\uses{def:physics:phyx-mini-0818:phyxminiproblems-problemphyxmini0818-lengthquantity, def:physics:phyx-mini-0818:phyxminiproblems-problemphyxmini0818-lengthinmeters, def:physics:phyx-mini-0818:phyxminiproblems-problemphyxmini0818-propellerdesignsetup, def:physics:phyx-mini-0818:phyxminiproblems-problemphyxmini0818-matchesprimarypropellerfigure, def:physics:phyx-mini-0818:phyxminiproblems-problemphyxmini0818-matchespropellerproblemdata, def:physics:phyx-mini-0818:phyxminiproblems-problemphyxmini0818-hasphysicalpropellerparameters, def:physics:phyx-mini-0818:phyxminiproblems-problemphyxmini0818-satisfiespropellerangularrateconversion, def:physics:phyx-mini-0818:phyxminiproblems-problemphyxmini0818-satisfiespropellertipkinematics, def:physics:phyx-mini-0818:phyxminiproblems-problemphyxmini0818-ismaximumadmissibleradius, def:physics:phyx-mini-0818:phyxminiproblems-problemphyxmini0818-radiusfromsaturatedspeedbudgetmeters, def:physics:phyx-mini-0818:phyxminiproblems-problemphyxmini0818-displayedradiusmeters, def:physics:phyx-mini-0818:phyxminiproblems-problemphyxmini0818-isclosestdisplayedradius}
The assigned autoformalize agent should translate this physics problem into Lean declarations in the covered file, with theorem and lemma proof bodies written as `by sorry`.
\end{theorem}
\begin{proof}
This is an autoformalization task, not a proof task. Produce statements that can later be proved by the physics prover without weakening the physical model.
\end{proof}

% --- Archon declaration topology begin ---
\section{Declaration topology}
\begin{definition}[Length Quantity]
  \label{def:physics:phyx-mini-0818:phyxminiproblems-problemphyxmini0818-lengthquantity}
  \lean{PhyXMiniProblems.ProblemPhyXMini0818.LengthQuantity}
  A nonnegative physical length, used for a candidate propeller radius
\end{definition}
\begin{proof}
  This is the declared model definition of Length Quantity.
\end{proof}
\begin{definition}[Revolution Rate Quantity]
  \label{def:physics:phyx-mini-0818:phyxminiproblems-problemphyxmini0818-revolutionratequantity}
  \lean{PhyXMiniProblems.ProblemPhyXMini0818.RevolutionRateQuantity}
  A nonnegative physical revolution rate, with dimension inverse time
\end{definition}
\begin{proof}
  This is the declared model definition of Revolution Rate Quantity.
\end{proof}
\begin{definition}[Angular Speed Quantity]
  \label{def:physics:phyx-mini-0818:phyxminiproblems-problemphyxmini0818-angularspeedquantity}
  \lean{PhyXMiniProblems.ProblemPhyXMini0818.AngularSpeedQuantity}
  A nonnegative angular-speed magnitude; radians are dimensionless
\end{definition}
\begin{proof}
  This is the declared model definition of Angular Speed Quantity.
\end{proof}
\begin{definition}[Speed Quantity]
  \label{def:physics:phyx-mini-0818:phyxminiproblems-problemphyxmini0818-speedquantity}
  \lean{PhyXMiniProblems.ProblemPhyXMini0818.SpeedQuantity}
  A nonnegative physical speed magnitude
\end{definition}
\begin{proof}
  This is the declared model definition of Speed Quantity.
\end{proof}
\begin{definition}[Spatial Velocity Quantity]
  \label{def:physics:phyx-mini-0818:phyxminiproblems-problemphyxmini0818-spatialvelocityquantity}
  \lean{PhyXMiniProblems.ProblemPhyXMini0818.SpatialVelocityQuantity}
  A dimensionful spatial velocity with three signed Cartesian components
\end{definition}
\begin{proof}
  This is the declared model definition of Spatial Velocity Quantity.
\end{proof}
\begin{definition}[length Readout]
  \label{def:physics:phyx-mini-0818:phyxminiproblems-problemphyxmini0818-lengthreadout}
  \lean{PhyXMiniProblems.ProblemPhyXMini0818.lengthReadout}
  \uses{def:physics:phyx-mini-0818:phyxminiproblems-problemphyxmini0818-lengthquantity}
  Read a nonnegative physical length in a coherent choice of units
\end{definition}
\begin{proof}
  This is the declared model definition of length Readout.
\end{proof}
\begin{definition}[revolution Rate Readout]
  \label{def:physics:phyx-mini-0818:phyxminiproblems-problemphyxmini0818-revolutionratereadout}
  \lean{PhyXMiniProblems.ProblemPhyXMini0818.revolutionRateReadout}
  \uses{def:physics:phyx-mini-0818:phyxminiproblems-problemphyxmini0818-revolutionratequantity}
  Read a revolution rate in inverse units of the selected time unit
\end{definition}
\begin{proof}
  This is the declared model definition of revolution Rate Readout.
\end{proof}
\begin{definition}[angular Speed Readout]
  \label{def:physics:phyx-mini-0818:phyxminiproblems-problemphyxmini0818-angularspeedreadout}
  \lean{PhyXMiniProblems.ProblemPhyXMini0818.angularSpeedReadout}
  \uses{def:physics:phyx-mini-0818:phyxminiproblems-problemphyxmini0818-angularspeedquantity}
  Read an angular speed in inverse units of the selected time unit
\end{definition}
\begin{proof}
  This is the declared model definition of angular Speed Readout.
\end{proof}
\begin{definition}[speed Readout]
  \label{def:physics:phyx-mini-0818:phyxminiproblems-problemphyxmini0818-speedreadout}
  \lean{PhyXMiniProblems.ProblemPhyXMini0818.speedReadout}
  \uses{def:physics:phyx-mini-0818:phyxminiproblems-problemphyxmini0818-speedquantity}
  Read a nonnegative speed magnitude in coherent units
\end{definition}
\begin{proof}
  This is the declared model definition of speed Readout.
\end{proof}
\begin{definition}[spatial Velocity Readout]
  \label{def:physics:phyx-mini-0818:phyxminiproblems-problemphyxmini0818-spatialvelocityreadout}
  \lean{PhyXMiniProblems.ProblemPhyXMini0818.spatialVelocityReadout}
  \uses{def:physics:phyx-mini-0818:phyxminiproblems-problemphyxmini0818-spatialvelocityquantity}
  Read a spatial velocity vector in coherent units
\end{definition}
\begin{proof}
  This is the declared model definition of spatial Velocity Readout.
\end{proof}
\begin{definition}[length In Meters]
  \label{def:physics:phyx-mini-0818:phyxminiproblems-problemphyxmini0818-lengthinmeters}
  \lean{PhyXMiniProblems.ProblemPhyXMini0818.lengthInMeters}
  \uses{def:physics:phyx-mini-0818:phyxminiproblems-problemphyxmini0818-lengthquantity, def:physics:phyx-mini-0818:phyxminiproblems-problemphyxmini0818-lengthreadout}
  SI metre readout of a physical length
\end{definition}
\begin{proof}
  This is the declared model definition of length In Meters.
\end{proof}
\begin{definition}[revolution Rate In Hertz]
  \label{def:physics:phyx-mini-0818:phyxminiproblems-problemphyxmini0818-revolutionrateinhertz}
  \lean{PhyXMiniProblems.ProblemPhyXMini0818.revolutionRateInHertz}
  \uses{def:physics:phyx-mini-0818:phyxminiproblems-problemphyxmini0818-revolutionratequantity, def:physics:phyx-mini-0818:phyxminiproblems-problemphyxmini0818-revolutionratereadout}
  SI revolutions-per-second readout of a physical revolution rate
\end{definition}
\begin{proof}
  This is the declared model definition of revolution Rate In Hertz.
\end{proof}
\begin{definition}[angular Speed In Radians Per Second]
  \label{def:physics:phyx-mini-0818:phyxminiproblems-problemphyxmini0818-angularspeedinradianspersecond}
  \lean{PhyXMiniProblems.ProblemPhyXMini0818.angularSpeedInRadiansPerSecond}
  \uses{def:physics:phyx-mini-0818:phyxminiproblems-problemphyxmini0818-angularspeedquantity, def:physics:phyx-mini-0818:phyxminiproblems-problemphyxmini0818-angularspeedreadout}
  SI radians-per-second readout of an angular-speed magnitude
\end{definition}
\begin{proof}
  This is the declared model definition of angular Speed In Radians Per Second.
\end{proof}
\begin{definition}[speed In Meters Per Second]
  \label{def:physics:phyx-mini-0818:phyxminiproblems-problemphyxmini0818-speedinmeterspersecond}
  \lean{PhyXMiniProblems.ProblemPhyXMini0818.speedInMetersPerSecond}
  \uses{def:physics:phyx-mini-0818:phyxminiproblems-problemphyxmini0818-speedquantity, def:physics:phyx-mini-0818:phyxminiproblems-problemphyxmini0818-speedreadout}
  SI metre-per-second readout of a speed magnitude
\end{definition}
\begin{proof}
  This is the declared model definition of speed In Meters Per Second.
\end{proof}
\begin{definition}[spatial Velocity In Meters Per Second]
  \label{def:physics:phyx-mini-0818:phyxminiproblems-problemphyxmini0818-spatialvelocityinmeterspersecond}
  \lean{PhyXMiniProblems.ProblemPhyXMini0818.spatialVelocityInMetersPerSecond}
  \uses{def:physics:phyx-mini-0818:phyxminiproblems-problemphyxmini0818-spatialvelocityquantity, def:physics:phyx-mini-0818:phyxminiproblems-problemphyxmini0818-spatialvelocityreadout}
  SI metre-per-second readout of a spatial velocity vector
\end{definition}
\begin{proof}
  This is the declared model definition of spatial Velocity In Meters Per Second.
\end{proof}
\begin{definition}[Propeller Location]
  \label{def:physics:phyx-mini-0818:phyxminiproblems-problemphyxmini0818-propellerlocation}
  \lean{PhyXMiniProblems.ProblemPhyXMini0818.PropellerLocation}
  Location of the propeller on the depicted airplane
\end{definition}
\begin{proof}
  This is the declared model definition of Propeller Location.
\end{proof}
\begin{definition}[Figure Arrow Direction]
  \label{def:physics:phyx-mini-0818:phyxminiproblems-problemphyxmini0818-figurearrowdirection}
  \lean{PhyXMiniProblems.ProblemPhyXMini0818.FigureArrowDirection}
  Directions of the two green arrows visible in the supplied image
\end{definition}
\begin{proof}
  This is the declared model definition of Figure Arrow Direction.
\end{proof}
\begin{definition}[Figure Label]
  \label{def:physics:phyx-mini-0818:phyxminiproblems-problemphyxmini0818-figurelabel}
  \lean{PhyXMiniProblems.ProblemPhyXMini0818.FigureLabel}
  Literal symbolic and numerical labels visible in the supplied image
\end{definition}
\begin{proof}
  This is the declared model definition of Figure Label.
\end{proof}
\begin{definition}[Propeller Figure]
  \label{def:physics:phyx-mini-0818:phyxminiproblems-problemphyxmini0818-propellerfigure}
  \lean{PhyXMiniProblems.ProblemPhyXMini0818.PropellerFigure}
  \uses{def:physics:phyx-mini-0818:phyxminiproblems-problemphyxmini0818-propellerlocation, def:physics:phyx-mini-0818:phyxminiproblems-problemphyxmini0818-figurearrowdirection, def:physics:phyx-mini-0818:phyxminiproblems-problemphyxmini0818-figurelabel}
  Qualitative and numerical information transcribed from primary image 818.png. No metric radius is inferred from the drawing: the radius arrow carries only the symbolic label r
\end{definition}
\begin{proof}
  This is the declared model definition of Propeller Figure.
\end{proof}
\begin{definition}[Propeller Design Setup]
  \label{def:physics:phyx-mini-0818:phyxminiproblems-problemphyxmini0818-propellerdesignsetup}
  \lean{PhyXMiniProblems.ProblemPhyXMini0818.PropellerDesignSetup}
  \uses{def:physics:phyx-mini-0818:phyxminiproblems-problemphyxmini0818-lengthquantity, def:physics:phyx-mini-0818:phyxminiproblems-problemphyxmini0818-revolutionratequantity, def:physics:phyx-mini-0818:phyxminiproblems-problemphyxmini0818-angularspeedquantity, def:physics:phyx-mini-0818:phyxminiproblems-problemphyxmini0818-speedquantity, def:physics:phyx-mini-0818:phyxminiproblems-problemphyxmini0818-spatialvelocityquantity, def:physics:phyx-mini-0818:phyxminiproblems-problemphyxmini0818-propellerfigure}
  Independent physical data for the design problem. The functions assigning tangential and resultant tip velocities to a candidate physical radius are not defined from the requested maximum. Their relations are imposed only by the governing-law interfaces below
\end{definition}
\begin{proof}
  This is the declared model definition of Propeller Design Setup.
\end{proof}
\begin{definition}[Matches Primary Propeller Figure]
  \label{def:physics:phyx-mini-0818:phyxminiproblems-problemphyxmini0818-matchesprimarypropellerfigure}
  \lean{PhyXMiniProblems.ProblemPhyXMini0818.MatchesPrimaryPropellerFigure}
  \uses{def:physics:phyx-mini-0818:phyxminiproblems-problemphyxmini0818-propellerdesignsetup}
  Exact transcription of the labels, placements, and arrow directions
\end{definition}
\begin{proof}
  This is the declared model definition of Matches Primary Propeller Figure.
\end{proof}
\begin{definition}[Matches Propeller Problem Data]
  \label{def:physics:phyx-mini-0818:phyxminiproblems-problemphyxmini0818-matchespropellerproblemdata}
  \lean{PhyXMiniProblems.ProblemPhyXMini0818.MatchesPropellerProblemData}
  \uses{def:physics:phyx-mini-0818:phyxminiproblems-problemphyxmini0818-revolutionrateinhertz, def:physics:phyx-mini-0818:phyxminiproblems-problemphyxmini0818-speedinmeterspersecond, def:physics:phyx-mini-0818:phyxminiproblems-problemphyxmini0818-propellerdesignsetup}
  Numerical data from the prose and its connection to the image. The printed 2400 rev/min is 2400/60 = 40 revolutions per coherent-SI second. The approximately-eighty-percent remark is retained as contextual data with an explicit, setup-supplied tolerance; it plays no role in solving for radius. No radius or answer choice occurs here
\end{definition}
\begin{proof}
  This is the declared model definition of Matches Propeller Problem Data.
\end{proof}
\begin{definition}[Has Physical Propeller Parameters]
  \label{def:physics:phyx-mini-0818:phyxminiproblems-problemphyxmini0818-hasphysicalpropellerparameters}
  \lean{PhyXMiniProblems.ProblemPhyXMini0818.HasPhysicalPropellerParameters}
  \uses{def:physics:phyx-mini-0818:phyxminiproblems-problemphyxmini0818-lengthinmeters, def:physics:phyx-mini-0818:phyxminiproblems-problemphyxmini0818-revolutionrateinhertz, def:physics:phyx-mini-0818:phyxminiproblems-problemphyxmini0818-angularspeedinradianspersecond, def:physics:phyx-mini-0818:phyxminiproblems-problemphyxmini0818-speedinmeterspersecond, def:physics:phyx-mini-0818:phyxminiproblems-problemphyxmini0818-propellerdesignsetup}
  Positivity and nondegeneracy conditions selecting the physical branch
\end{definition}
\begin{proof}
  This is the declared model definition of Has Physical Propeller Parameters.
\end{proof}
\begin{definition}[Satisfies Propeller Angular Rate Conversion]
  \label{def:physics:phyx-mini-0818:phyxminiproblems-problemphyxmini0818-satisfiespropellerangularrateconversion}
  \lean{PhyXMiniProblems.ProblemPhyXMini0818.SatisfiesPropellerAngularRateConversion}
  \uses{def:physics:phyx-mini-0818:phyxminiproblems-problemphyxmini0818-revolutionratereadout, def:physics:phyx-mini-0818:phyxminiproblems-problemphyxmini0818-angularspeedreadout, def:physics:phyx-mini-0818:phyxminiproblems-problemphyxmini0818-propellerdesignsetup}
  One complete propeller revolution is 2*pi radians
\end{definition}
\begin{proof}
  This is the declared model definition of Satisfies Propeller Angular Rate Conversion.
\end{proof}
\begin{definition}[Satisfies Propeller Tip Kinematics]
  \label{def:physics:phyx-mini-0818:phyxminiproblems-problemphyxmini0818-satisfiespropellertipkinematics}
  \lean{PhyXMiniProblems.ProblemPhyXMini0818.SatisfiesPropellerTipKinematics}
  \uses{def:physics:phyx-mini-0818:phyxminiproblems-problemphyxmini0818-lengthquantity, def:physics:phyx-mini-0818:phyxminiproblems-problemphyxmini0818-lengthreadout, def:physics:phyx-mini-0818:phyxminiproblems-problemphyxmini0818-angularspeedreadout, def:physics:phyx-mini-0818:phyxminiproblems-problemphyxmini0818-speedreadout, def:physics:phyx-mini-0818:phyxminiproblems-problemphyxmini0818-spatialvelocityreadout, def:physics:phyx-mini-0818:phyxminiproblems-problemphyxmini0818-speedinmeterspersecond, def:physics:phyx-mini-0818:phyxminiproblems-problemphyxmini0818-spatialvelocityinmeterspersecond, def:physics:phyx-mini-0818:phyxminiproblems-problemphyxmini0818-propellerdesignsetup}
  Vector kinematics for each candidate radius. The forward velocity is the airplane velocity through the air, the tangential velocity is relative to the airframe, and their sum is the blade-tip velocity through the air. The perpendicularity field is the geometric fact displayed by the horizontal and vertical green arrows. It is the hypothesis needed for Mathlib's norm\_add\_sq\_eq\_norm\_sq\_add\_norm\_sq\_of\_inner\_eq\_zero Pythagorean theorem.
\end{definition}
\begin{proof}
  This is the declared model definition of Satisfies Propeller Tip Kinematics.
\end{proof}
\begin{lemma}[angular Speed is eighty pi radians per second]
  \label{lem:physics:phyx-mini-0818:phyxminiproblems-problemphyxmini0818-angularspeed-is-eighty-pi-radians-per-second}
  \lean{PhyXMiniProblems.ProblemPhyXMini0818.angularSpeed_is_eighty_pi_radians_per_second}
  \uses{def:physics:phyx-mini-0818:phyxminiproblems-problemphyxmini0818-angularspeedinradianspersecond, def:physics:phyx-mini-0818:phyxminiproblems-problemphyxmini0818-propellerdesignsetup, def:physics:phyx-mini-0818:phyxminiproblems-problemphyxmini0818-matchesprimarypropellerfigure, def:physics:phyx-mini-0818:phyxminiproblems-problemphyxmini0818-matchespropellerproblemdata, def:physics:phyx-mini-0818:phyxminiproblems-problemphyxmini0818-satisfiespropellerangularrateconversion}
  The angular rate conversion and the printed 2400 rev/min imply omega = 80*pi rad/s. This is a derived result, not problem data
\end{lemma}
\begin{proof}
  The conclusion follows from the governing relations and figure readouts named in the statement of angular Speed is eighty pi radians per second.
\end{proof}
\begin{lemma}[tip Speed Squared pythagorean]
  \label{lem:physics:phyx-mini-0818:phyxminiproblems-problemphyxmini0818-tipspeedsquared-pythagorean}
  \lean{PhyXMiniProblems.ProblemPhyXMini0818.tipSpeedSquared_pythagorean}
  \uses{def:physics:phyx-mini-0818:phyxminiproblems-problemphyxmini0818-lengthquantity, def:physics:phyx-mini-0818:phyxminiproblems-problemphyxmini0818-lengthinmeters, def:physics:phyx-mini-0818:phyxminiproblems-problemphyxmini0818-angularspeedinradianspersecond, def:physics:phyx-mini-0818:phyxminiproblems-problemphyxmini0818-speedinmeterspersecond, def:physics:phyx-mini-0818:phyxminiproblems-problemphyxmini0818-spatialvelocityinmeterspersecond, def:physics:phyx-mini-0818:phyxminiproblems-problemphyxmini0818-propellerdesignsetup, def:physics:phyx-mini-0818:phyxminiproblems-problemphyxmini0818-satisfiespropellertipkinematics}
  Pythagoras for the forward and tangential components of the blade-tip velocity. This lemma is a generic consequence of the vector laws and does not solve the maximum-radius optimization
\end{lemma}
\begin{proof}
  The conclusion follows from the governing relations and figure readouts named in the statement of tip Speed Squared pythagorean.
\end{proof}
\begin{definition}[Is Maximum Admissible Radius]
  \label{def:physics:phyx-mini-0818:phyxminiproblems-problemphyxmini0818-ismaximumadmissibleradius}
  \lean{PhyXMiniProblems.ProblemPhyXMini0818.IsMaximumAdmissibleRadius}
  \uses{def:physics:phyx-mini-0818:phyxminiproblems-problemphyxmini0818-lengthquantity, def:physics:phyx-mini-0818:phyxminiproblems-problemphyxmini0818-lengthinmeters, def:physics:phyx-mini-0818:phyxminiproblems-problemphyxmini0818-propellerdesignsetup}
  A physical radius is maximal when it is admissible and bounds every other admissible physical radius in the SI metre readout
\end{definition}
\begin{proof}
  This is the declared model definition of Is Maximum Admissible Radius.
\end{proof}
\begin{definition}[radius From Saturated Speed Budget Meters]
  \label{def:physics:phyx-mini-0818:phyxminiproblems-problemphyxmini0818-radiusfromsaturatedspeedbudgetmeters}
  \lean{PhyXMiniProblems.ProblemPhyXMini0818.radiusFromSaturatedSpeedBudgetMeters}
  The SI radius candidate obtained algebraically by saturating the speed cap. Its maximality is not definitional; that is the substantive theorem below
\end{definition}
\begin{proof}
  This is the declared model definition of radius From Saturated Speed Budget Meters.
\end{proof}
\begin{definition}[Answer Choice]
  \label{def:physics:phyx-mini-0818:phyxminiproblems-problemphyxmini0818-answerchoice}
  \lean{PhyXMiniProblems.ProblemPhyXMini0818.AnswerChoice}
  Labels of the four radius choices displayed with the problem
\end{definition}
\begin{proof}
  This is the declared model definition of Answer Choice.
\end{proof}
\begin{definition}[displayed Radius Meters]
  \label{def:physics:phyx-mini-0818:phyxminiproblems-problemphyxmini0818-displayedradiusmeters}
  \lean{PhyXMiniProblems.ProblemPhyXMini0818.displayedRadiusMeters}
  \uses{def:physics:phyx-mini-0818:phyxminiproblems-problemphyxmini0818-answerchoice}
  Metre readout printed beside each displayed answer choice
\end{definition}
\begin{proof}
  This is the declared model definition of displayed Radius Meters.
\end{proof}
\begin{definition}[recorded Dataset Answer]
  \label{def:physics:phyx-mini-0818:phyxminiproblems-problemphyxmini0818-recordeddatasetanswer}
  \lean{PhyXMiniProblems.ProblemPhyXMini0818.recordedDatasetAnswer}
  \uses{def:physics:phyx-mini-0818:phyxminiproblems-problemphyxmini0818-answerchoice}
  The answer label recorded in the supplied dataset metadata
\end{definition}
\begin{proof}
  This is the declared model definition of recorded Dataset Answer.
\end{proof}
\begin{definition}[Is Closest Displayed Radius]
  \label{def:physics:phyx-mini-0818:phyxminiproblems-problemphyxmini0818-isclosestdisplayedradius}
  \lean{PhyXMiniProblems.ProblemPhyXMini0818.IsClosestDisplayedRadius}
  \uses{def:physics:phyx-mini-0818:phyxminiproblems-problemphyxmini0818-answerchoice, def:physics:phyx-mini-0818:phyxminiproblems-problemphyxmini0818-displayedradiusmeters}
  A displayed radius is closest to a computed radius readout
\end{definition}
\begin{proof}
  This is the declared model definition of Is Closest Displayed Radius.
\end{proof}
% --- Archon declaration topology end ---
% --- Archon physics formalization source end ---
